

	\documentclass[10pt,parskip=half,
	toc=sectionentrywithdots,
	bibliography=totocnumbered,
	captions=tableheading,
    numbers=noendperiod,
    headings=standardclasses]{scrartcl}

    \usepackage[T1]{fontenc} % Nicer default font (+ math font) than Computer Modern for most use cases
    \usepackage{mathpazo}
    \usepackage{graphicx}
    \usepackage[skip=3pt]{caption}
    \usepackage{adjustbox} % Used to constrain images to a maximum size
    \usepackage[table]{xcolor} % Allow colors to be defined
    \usepackage{enumerate} % Needed for markdown enumerations to work
    \usepackage{amsmath} % Equations
    \usepackage{amssymb} % Equations
    \usepackage{textcomp} % defines textquotesingle
    % Hack from http://tex.stackexchange.com/a/47451/13684:
    \AtBeginDocument{%
        \def\PYZsq{\textquotesingle}% Upright quotes in Pygmentized code
    }
    \usepackage{upquote} % Upright quotes for verbatim code
    \usepackage{eurosym} % defines \euro
    \usepackage[mathletters]{ucs} % Extended unicode (utf-8) support
    \usepackage[utf8x]{inputenc} % Allow utf-8 characters in the tex document
    \usepackage{fancyvrb} % verbatim replacement that allows latex
    \usepackage{grffile} % extends the file name processing of package graphics
                         % to support a larger range
    % The hyperref package gives us a pdf with properly built
    % internal navigation ('pdf bookmarks' for the table of contents,
    % internal cross-reference links, web links for URLs, etc.)
    \usepackage{hyperref}
    \usepackage{longtable} % longtable support required by pandoc >1.10
    \usepackage{booktabs}  % table support for pandoc > 1.12.2
    \usepackage[inline]{enumitem} % IRkernel/repr support (it uses the enumerate* environment)
    \usepackage[normalem]{ulem} % ulem is needed to support strikethroughs (\sout)
                                % normalem makes italics be italics, not underlines

    \usepackage{translations}
	\usepackage{microtype} % improves the spacing between words and letters
	\usepackage{placeins} % placement of figures
    % could use \usepackage[section]{placeins} but placing in subsection in command section
	% Places the float at precisely the location in the LaTeX code (with H)
	\usepackage{float}
	\usepackage[colorinlistoftodos,obeyFinal,textwidth=.8in]{todonotes} % to mark to-dos
	% number figures, tables and equations by section
	\usepackage{chngcntr}
	% header/footer
	\usepackage[footsepline=0.25pt, headsepline=0.25pt]{scrlayer-scrpage}

	% bibliography formatting
	\usepackage[natbibapa, nodoi]{apacite} % WORKING -- suppresses doi
	% \usepackage[numbers, square, super, sort&compress]{natbib} % ORIG
	% hyperlink doi's
	\usepackage{doi}

    % define a code float
    \usepackage{newfloat} % to define a new float types
    \DeclareFloatingEnvironment[
        fileext=frm,placement={!ht},
        within=section,name=Code]{codecell}
    \DeclareFloatingEnvironment[
        fileext=frm,placement={!ht},
        within=section,name=Text]{textcell}
    \DeclareFloatingEnvironment[
        fileext=frm,placement={!ht},
        within=section,name=Text]{errorcell}

    \usepackage{listings} % a package for wrapping code in a box
    \usepackage[framemethod=tikz]{mdframed} % to fram code

% Pygments definitions

\makeatletter
\def\PY@reset{\let\PY@it=\relax \let\PY@bf=\relax%
    \let\PY@ul=\relax \let\PY@tc=\relax%
    \let\PY@bc=\relax \let\PY@ff=\relax}
\def\PY@tok#1{\csname PY@tok@#1\endcsname}
\def\PY@toks#1+{\ifx\relax#1\empty\else%
    \PY@tok{#1}\expandafter\PY@toks\fi}
\def\PY@do#1{\PY@bc{\PY@tc{\PY@ul{%
    \PY@it{\PY@bf{\PY@ff{#1}}}}}}}
\def\PY#1#2{\PY@reset\PY@toks#1+\relax+\PY@do{#2}}

\expandafter\def\csname PY@tok@w\endcsname{\def\PY@tc##1{\textcolor[rgb]{0.73,0.73,0.73}{##1}}}
\expandafter\def\csname PY@tok@c\endcsname{\let\PY@it=\textit\def\PY@tc##1{\textcolor[rgb]{0.25,0.50,0.50}{##1}}}
\expandafter\def\csname PY@tok@cp\endcsname{\def\PY@tc##1{\textcolor[rgb]{0.74,0.48,0.00}{##1}}}
\expandafter\def\csname PY@tok@k\endcsname{\let\PY@bf=\textbf\def\PY@tc##1{\textcolor[rgb]{0.00,0.50,0.00}{##1}}}
\expandafter\def\csname PY@tok@kp\endcsname{\def\PY@tc##1{\textcolor[rgb]{0.00,0.50,0.00}{##1}}}
\expandafter\def\csname PY@tok@kt\endcsname{\def\PY@tc##1{\textcolor[rgb]{0.69,0.00,0.25}{##1}}}
\expandafter\def\csname PY@tok@o\endcsname{\def\PY@tc##1{\textcolor[rgb]{0.40,0.40,0.40}{##1}}}
\expandafter\def\csname PY@tok@ow\endcsname{\let\PY@bf=\textbf\def\PY@tc##1{\textcolor[rgb]{0.67,0.13,1.00}{##1}}}
\expandafter\def\csname PY@tok@nb\endcsname{\def\PY@tc##1{\textcolor[rgb]{0.00,0.50,0.00}{##1}}}
\expandafter\def\csname PY@tok@nf\endcsname{\def\PY@tc##1{\textcolor[rgb]{0.00,0.00,1.00}{##1}}}
\expandafter\def\csname PY@tok@nc\endcsname{\let\PY@bf=\textbf\def\PY@tc##1{\textcolor[rgb]{0.00,0.00,1.00}{##1}}}
\expandafter\def\csname PY@tok@nn\endcsname{\let\PY@bf=\textbf\def\PY@tc##1{\textcolor[rgb]{0.00,0.00,1.00}{##1}}}
\expandafter\def\csname PY@tok@ne\endcsname{\let\PY@bf=\textbf\def\PY@tc##1{\textcolor[rgb]{0.82,0.25,0.23}{##1}}}
\expandafter\def\csname PY@tok@nv\endcsname{\def\PY@tc##1{\textcolor[rgb]{0.10,0.09,0.49}{##1}}}
\expandafter\def\csname PY@tok@no\endcsname{\def\PY@tc##1{\textcolor[rgb]{0.53,0.00,0.00}{##1}}}
\expandafter\def\csname PY@tok@nl\endcsname{\def\PY@tc##1{\textcolor[rgb]{0.63,0.63,0.00}{##1}}}
\expandafter\def\csname PY@tok@ni\endcsname{\let\PY@bf=\textbf\def\PY@tc##1{\textcolor[rgb]{0.60,0.60,0.60}{##1}}}
\expandafter\def\csname PY@tok@na\endcsname{\def\PY@tc##1{\textcolor[rgb]{0.49,0.56,0.16}{##1}}}
\expandafter\def\csname PY@tok@nt\endcsname{\let\PY@bf=\textbf\def\PY@tc##1{\textcolor[rgb]{0.00,0.50,0.00}{##1}}}
\expandafter\def\csname PY@tok@nd\endcsname{\def\PY@tc##1{\textcolor[rgb]{0.67,0.13,1.00}{##1}}}
\expandafter\def\csname PY@tok@s\endcsname{\def\PY@tc##1{\textcolor[rgb]{0.73,0.13,0.13}{##1}}}
\expandafter\def\csname PY@tok@sd\endcsname{\let\PY@it=\textit\def\PY@tc##1{\textcolor[rgb]{0.73,0.13,0.13}{##1}}}
\expandafter\def\csname PY@tok@si\endcsname{\let\PY@bf=\textbf\def\PY@tc##1{\textcolor[rgb]{0.73,0.40,0.53}{##1}}}
\expandafter\def\csname PY@tok@se\endcsname{\let\PY@bf=\textbf\def\PY@tc##1{\textcolor[rgb]{0.73,0.40,0.13}{##1}}}
\expandafter\def\csname PY@tok@sr\endcsname{\def\PY@tc##1{\textcolor[rgb]{0.73,0.40,0.53}{##1}}}
\expandafter\def\csname PY@tok@ss\endcsname{\def\PY@tc##1{\textcolor[rgb]{0.10,0.09,0.49}{##1}}}
\expandafter\def\csname PY@tok@sx\endcsname{\def\PY@tc##1{\textcolor[rgb]{0.00,0.50,0.00}{##1}}}
\expandafter\def\csname PY@tok@m\endcsname{\def\PY@tc##1{\textcolor[rgb]{0.40,0.40,0.40}{##1}}}
\expandafter\def\csname PY@tok@gh\endcsname{\let\PY@bf=\textbf\def\PY@tc##1{\textcolor[rgb]{0.00,0.00,0.50}{##1}}}
\expandafter\def\csname PY@tok@gu\endcsname{\let\PY@bf=\textbf\def\PY@tc##1{\textcolor[rgb]{0.50,0.00,0.50}{##1}}}
\expandafter\def\csname PY@tok@gd\endcsname{\def\PY@tc##1{\textcolor[rgb]{0.63,0.00,0.00}{##1}}}
\expandafter\def\csname PY@tok@gi\endcsname{\def\PY@tc##1{\textcolor[rgb]{0.00,0.63,0.00}{##1}}}
\expandafter\def\csname PY@tok@gr\endcsname{\def\PY@tc##1{\textcolor[rgb]{1.00,0.00,0.00}{##1}}}
\expandafter\def\csname PY@tok@ge\endcsname{\let\PY@it=\textit}
\expandafter\def\csname PY@tok@gs\endcsname{\let\PY@bf=\textbf}
\expandafter\def\csname PY@tok@gp\endcsname{\let\PY@bf=\textbf\def\PY@tc##1{\textcolor[rgb]{0.00,0.00,0.50}{##1}}}
\expandafter\def\csname PY@tok@go\endcsname{\def\PY@tc##1{\textcolor[rgb]{0.53,0.53,0.53}{##1}}}
\expandafter\def\csname PY@tok@gt\endcsname{\def\PY@tc##1{\textcolor[rgb]{0.00,0.27,0.87}{##1}}}
\expandafter\def\csname PY@tok@err\endcsname{\def\PY@bc##1{\setlength{\fboxsep}{0pt}\fcolorbox[rgb]{1.00,0.00,0.00}{1,1,1}{\strut ##1}}}
\expandafter\def\csname PY@tok@kc\endcsname{\let\PY@bf=\textbf\def\PY@tc##1{\textcolor[rgb]{0.00,0.50,0.00}{##1}}}
\expandafter\def\csname PY@tok@kd\endcsname{\let\PY@bf=\textbf\def\PY@tc##1{\textcolor[rgb]{0.00,0.50,0.00}{##1}}}
\expandafter\def\csname PY@tok@kn\endcsname{\let\PY@bf=\textbf\def\PY@tc##1{\textcolor[rgb]{0.00,0.50,0.00}{##1}}}
\expandafter\def\csname PY@tok@kr\endcsname{\let\PY@bf=\textbf\def\PY@tc##1{\textcolor[rgb]{0.00,0.50,0.00}{##1}}}
\expandafter\def\csname PY@tok@bp\endcsname{\def\PY@tc##1{\textcolor[rgb]{0.00,0.50,0.00}{##1}}}
\expandafter\def\csname PY@tok@fm\endcsname{\def\PY@tc##1{\textcolor[rgb]{0.00,0.00,1.00}{##1}}}
\expandafter\def\csname PY@tok@vc\endcsname{\def\PY@tc##1{\textcolor[rgb]{0.10,0.09,0.49}{##1}}}
\expandafter\def\csname PY@tok@vg\endcsname{\def\PY@tc##1{\textcolor[rgb]{0.10,0.09,0.49}{##1}}}
\expandafter\def\csname PY@tok@vi\endcsname{\def\PY@tc##1{\textcolor[rgb]{0.10,0.09,0.49}{##1}}}
\expandafter\def\csname PY@tok@vm\endcsname{\def\PY@tc##1{\textcolor[rgb]{0.10,0.09,0.49}{##1}}}
\expandafter\def\csname PY@tok@sa\endcsname{\def\PY@tc##1{\textcolor[rgb]{0.73,0.13,0.13}{##1}}}
\expandafter\def\csname PY@tok@sb\endcsname{\def\PY@tc##1{\textcolor[rgb]{0.73,0.13,0.13}{##1}}}
\expandafter\def\csname PY@tok@sc\endcsname{\def\PY@tc##1{\textcolor[rgb]{0.73,0.13,0.13}{##1}}}
\expandafter\def\csname PY@tok@dl\endcsname{\def\PY@tc##1{\textcolor[rgb]{0.73,0.13,0.13}{##1}}}
\expandafter\def\csname PY@tok@s2\endcsname{\def\PY@tc##1{\textcolor[rgb]{0.73,0.13,0.13}{##1}}}
\expandafter\def\csname PY@tok@sh\endcsname{\def\PY@tc##1{\textcolor[rgb]{0.73,0.13,0.13}{##1}}}
\expandafter\def\csname PY@tok@s1\endcsname{\def\PY@tc##1{\textcolor[rgb]{0.73,0.13,0.13}{##1}}}
\expandafter\def\csname PY@tok@mb\endcsname{\def\PY@tc##1{\textcolor[rgb]{0.40,0.40,0.40}{##1}}}
\expandafter\def\csname PY@tok@mf\endcsname{\def\PY@tc##1{\textcolor[rgb]{0.40,0.40,0.40}{##1}}}
\expandafter\def\csname PY@tok@mh\endcsname{\def\PY@tc##1{\textcolor[rgb]{0.40,0.40,0.40}{##1}}}
\expandafter\def\csname PY@tok@mi\endcsname{\def\PY@tc##1{\textcolor[rgb]{0.40,0.40,0.40}{##1}}}
\expandafter\def\csname PY@tok@il\endcsname{\def\PY@tc##1{\textcolor[rgb]{0.40,0.40,0.40}{##1}}}
\expandafter\def\csname PY@tok@mo\endcsname{\def\PY@tc##1{\textcolor[rgb]{0.40,0.40,0.40}{##1}}}
\expandafter\def\csname PY@tok@ch\endcsname{\let\PY@it=\textit\def\PY@tc##1{\textcolor[rgb]{0.25,0.50,0.50}{##1}}}
\expandafter\def\csname PY@tok@cm\endcsname{\let\PY@it=\textit\def\PY@tc##1{\textcolor[rgb]{0.25,0.50,0.50}{##1}}}
\expandafter\def\csname PY@tok@cpf\endcsname{\let\PY@it=\textit\def\PY@tc##1{\textcolor[rgb]{0.25,0.50,0.50}{##1}}}
\expandafter\def\csname PY@tok@c1\endcsname{\let\PY@it=\textit\def\PY@tc##1{\textcolor[rgb]{0.25,0.50,0.50}{##1}}}
\expandafter\def\csname PY@tok@cs\endcsname{\let\PY@it=\textit\def\PY@tc##1{\textcolor[rgb]{0.25,0.50,0.50}{##1}}}

\def\PYZbs{\char`\\}
\def\PYZus{\char`\_}
\def\PYZob{\char`\{}
\def\PYZcb{\char`\}}
\def\PYZca{\char`\^}
\def\PYZam{\char`\&}
\def\PYZlt{\char`\<}
\def\PYZgt{\char`\>}
\def\PYZsh{\char`\#}
\def\PYZpc{\char`\%}
\def\PYZdl{\char`\$}
\def\PYZhy{\char`\-}
\def\PYZsq{\char`\'}
\def\PYZdq{\char`\"}
\def\PYZti{\char`\~}
% for compatibility with earlier versions
\def\PYZat{@}
\def\PYZlb{[}
\def\PYZrb{]}
\makeatother

% ANSI colors
\definecolor{ansi-black}{HTML}{3E424D}
\definecolor{ansi-black-intense}{HTML}{282C36}
\definecolor{ansi-red}{HTML}{E75C58}
\definecolor{ansi-red-intense}{HTML}{B22B31}
\definecolor{ansi-green}{HTML}{00A250}
\definecolor{ansi-green-intense}{HTML}{007427}
\definecolor{ansi-yellow}{HTML}{DDB62B}
\definecolor{ansi-yellow-intense}{HTML}{B27D12}
\definecolor{ansi-blue}{HTML}{208FFB}
\definecolor{ansi-blue-intense}{HTML}{0065CA}
\definecolor{ansi-magenta}{HTML}{D160C4}
\definecolor{ansi-magenta-intense}{HTML}{A03196}
\definecolor{ansi-cyan}{HTML}{60C6C8}
\definecolor{ansi-cyan-intense}{HTML}{258F8F}
\definecolor{ansi-white}{HTML}{C5C1B4}
\definecolor{ansi-white-intense}{HTML}{A1A6B2}

% commands and environments needed by pandoc snippets
% extracted from the output of `pandoc -s`
\providecommand{\tightlist}{%
  \setlength{\itemsep}{0pt}\setlength{\parskip}{0pt}}
\DefineVerbatimEnvironment{Highlighting}{Verbatim}{commandchars=\\\{\}}
% Add ',fontsize=\small' for more characters per line
\newenvironment{Shaded}{}{}
\newcommand{\KeywordTok}[1]{\textcolor[rgb]{0.00,0.44,0.13}{\textbf{{#1}}}}
\newcommand{\DataTypeTok}[1]{\textcolor[rgb]{0.56,0.13,0.00}{{#1}}}
\newcommand{\DecValTok}[1]{\textcolor[rgb]{0.25,0.63,0.44}{{#1}}}
\newcommand{\BaseNTok}[1]{\textcolor[rgb]{0.25,0.63,0.44}{{#1}}}
\newcommand{\FloatTok}[1]{\textcolor[rgb]{0.25,0.63,0.44}{{#1}}}
\newcommand{\CharTok}[1]{\textcolor[rgb]{0.25,0.44,0.63}{{#1}}}
\newcommand{\StringTok}[1]{\textcolor[rgb]{0.25,0.44,0.63}{{#1}}}
\newcommand{\CommentTok}[1]{\textcolor[rgb]{0.38,0.63,0.69}{\textit{{#1}}}}
\newcommand{\OtherTok}[1]{\textcolor[rgb]{0.00,0.44,0.13}{{#1}}}
\newcommand{\AlertTok}[1]{\textcolor[rgb]{1.00,0.00,0.00}{\textbf{{#1}}}}
\newcommand{\FunctionTok}[1]{\textcolor[rgb]{0.02,0.16,0.49}{{#1}}}
\newcommand{\RegionMarkerTok}[1]{{#1}}
\newcommand{\ErrorTok}[1]{\textcolor[rgb]{1.00,0.00,0.00}{\textbf{{#1}}}}
\newcommand{\NormalTok}[1]{{#1}}

% Additional commands for more recent versions of Pandoc
\newcommand{\ConstantTok}[1]{\textcolor[rgb]{0.53,0.00,0.00}{{#1}}}
\newcommand{\SpecialCharTok}[1]{\textcolor[rgb]{0.25,0.44,0.63}{{#1}}}
\newcommand{\VerbatimStringTok}[1]{\textcolor[rgb]{0.25,0.44,0.63}{{#1}}}
\newcommand{\SpecialStringTok}[1]{\textcolor[rgb]{0.73,0.40,0.53}{{#1}}}
\newcommand{\ImportTok}[1]{{#1}}
\newcommand{\DocumentationTok}[1]{\textcolor[rgb]{0.73,0.13,0.13}{\textit{{#1}}}}
\newcommand{\AnnotationTok}[1]{\textcolor[rgb]{0.38,0.63,0.69}{\textbf{\textit{{#1}}}}}
\newcommand{\CommentVarTok}[1]{\textcolor[rgb]{0.38,0.63,0.69}{\textbf{\textit{{#1}}}}}
\newcommand{\VariableTok}[1]{\textcolor[rgb]{0.10,0.09,0.49}{{#1}}}
\newcommand{\ControlFlowTok}[1]{\textcolor[rgb]{0.00,0.44,0.13}{\textbf{{#1}}}}
\newcommand{\OperatorTok}[1]{\textcolor[rgb]{0.40,0.40,0.40}{{#1}}}
\newcommand{\BuiltInTok}[1]{{#1}}
\newcommand{\ExtensionTok}[1]{{#1}}
\newcommand{\PreprocessorTok}[1]{\textcolor[rgb]{0.74,0.48,0.00}{{#1}}}
\newcommand{\AttributeTok}[1]{\textcolor[rgb]{0.49,0.56,0.16}{{#1}}}
\newcommand{\InformationTok}[1]{\textcolor[rgb]{0.38,0.63,0.69}{\textbf{\textit{{#1}}}}}
\newcommand{\WarningTok}[1]{\textcolor[rgb]{0.38,0.63,0.69}{\textbf{\textit{{#1}}}}}

% Define a nice break command that doesn't care if a line doesn't already
% exist.
\def\br{\hspace*{\fill} \\* }

% Math Jax compatability definitions
\def\gt{>}
\def\lt{<}

    \setcounter{secnumdepth}{5}

    % Colors for the hyperref package
    \definecolor{urlcolor}{rgb}{0,.145,.698}
    \definecolor{linkcolor}{rgb}{.71,0.21,0.01}
    \definecolor{citecolor}{rgb}{.12,.54,.11}

\DeclareTranslationFallback{Author}{Author}
\DeclareTranslation{Portuges}{Author}{Autor}

\DeclareTranslationFallback{List of Codes}{List of Codes}
\DeclareTranslation{Catalan}{List of Codes}{Llista de Codis}
\DeclareTranslation{Danish}{List of Codes}{Liste over Koder}
\DeclareTranslation{German}{List of Codes}{Liste der Codes}
\DeclareTranslation{Spanish}{List of Codes}{Lista de C\'{o}digos}
\DeclareTranslation{French}{List of Codes}{Liste des Codes}
\DeclareTranslation{Italian}{List of Codes}{Elenco dei Codici}
\DeclareTranslation{Dutch}{List of Codes}{Lijst van Codes}
\DeclareTranslation{Portuges}{List of Codes}{Lista de C\'{o}digos}

\DeclareTranslationFallback{Supervisors}{Supervisors}
\DeclareTranslation{Catalan}{Supervisors}{Supervisors}
\DeclareTranslation{Danish}{Supervisors}{Vejledere}
\DeclareTranslation{German}{Supervisors}{Vorgesetzten}
\DeclareTranslation{Spanish}{Supervisors}{Supervisores}
\DeclareTranslation{French}{Supervisors}{Superviseurs}
\DeclareTranslation{Italian}{Supervisors}{Le autorit\`{a} di vigilanza}
\DeclareTranslation{Dutch}{Supervisors}{supervisors}
\DeclareTranslation{Portuguese}{Supervisors}{Supervisores}

\definecolor{codegreen}{rgb}{0,0.6,0}
\definecolor{codegray}{rgb}{0.5,0.5,0.5}
\definecolor{codepurple}{rgb}{0.58,0,0.82}
\definecolor{backcolour}{rgb}{0.95,0.95,0.95}

\lstdefinestyle{mystyle}{
    commentstyle=\color{codegreen},
    keywordstyle=\color{magenta},
    numberstyle=\tiny\color{codegray},
    stringstyle=\color{codepurple},
    basicstyle=\ttfamily,
    breakatwhitespace=false,
    keepspaces=true,
    numbers=left,
    numbersep=10pt,
    showspaces=false,
    showstringspaces=false,
    showtabs=false,
    tabsize=2,
    breaklines=true,
    literate={\-}{}{0\discretionary{-}{}{-}},
  postbreak=\mbox{\textcolor{red}{$\hookrightarrow$}\space},
}

\lstset{style=mystyle}

\surroundwithmdframed[
  hidealllines=true,
  backgroundcolor=backcolour,
  innerleftmargin=0pt,
  innerrightmargin=0pt,
  innertopmargin=0pt,
  innerbottommargin=0pt]{lstlisting}

 % Used to adjust the document margins
\usepackage{geometry}
\geometry{tmargin=1in,bmargin=1in,lmargin=1in,rmargin=1in,
includehead,includefoot,footskip=25pt} % use nohead if... well... no heading
% you can use showframe option to check the margins visually

	% ensure new section starts on new page
	\addtokomafont{section}{\clearpage}

    % Prevent overflowing lines due to hard-to-break entities
    \sloppy

    % Setup hyperref package
    \hypersetup{
      breaklinks=true,  % so long urls are correctly broken across lines
      colorlinks=true,
      urlcolor=urlcolor,
      linkcolor=linkcolor,
      citecolor=citecolor,
      }

    % ensure figures are placed within subsections
    \makeatletter
    \AtBeginDocument{%
      \expandafter\renewcommand\expandafter\subsection\expandafter
        {\expandafter\@fb@secFB\subsection}%
      \newcommand\@fb@secFB{\FloatBarrier
        \gdef\@fb@afterHHook{\@fb@topbarrier \gdef\@fb@afterHHook{}}}%
      \g@addto@macro\@afterheading{\@fb@afterHHook}%
      \gdef\@fb@afterHHook{}%
    }
    \makeatother

	% number figures, tables and equations by section
	\usepackage{chngcntr}
	\counterwithout{figure}{section}
	\counterwithout{table}{section}
	\counterwithout{equation}{section}
	\makeatletter
	\@addtoreset{table}{section}
	\@addtoreset{figure}{section}
	\@addtoreset{equation}{section}
	\makeatother
    % XXX -- Use this for section numbered figures and tables
	\renewcommand\thetable{\thesection.\arabic{table}}
	\renewcommand\thefigure{\thesection.\arabic{figure}}
	\renewcommand\theequation{\thesection.\arabic{equation}}

    % XXX -- Use this for absolute numbered figures and tables.
    % \renewcommand\thetable{\arabic{table}}
	% \renewcommand\thefigure{\arabic{figure}}
	% \renewcommand\theequation{\arabic{equation}}

        % set global options for float placement
        \makeatletter
          \providecommand*\setfloatlocations[2]{\@namedef{fps@#1}{#2}}
        \makeatother

    % align captions to left (indented)
	\captionsetup{justification=raggedright,
	singlelinecheck=false,format=hang,labelfont={it,bf}}

	% shift footer down so space between separation line
	\ModifyLayer[addvoffset=.6ex]{scrheadings.foot.odd}
	\ModifyLayer[addvoffset=.6ex]{scrheadings.foot.even}
	\ModifyLayer[addvoffset=.6ex]{scrheadings.foot.oneside}
	\ModifyLayer[addvoffset=.6ex]{plain.scrheadings.foot.odd}
	\ModifyLayer[addvoffset=.6ex]{plain.scrheadings.foot.even}
	\ModifyLayer[addvoffset=.6ex]{plain.scrheadings.foot.oneside}
	\pagestyle{scrheadings}
	\clearscrheadfoot{}

    % Running head in header
        \ihead{PsiPyPublish}

    % Footer
	\ifoot{\leftmark}
	\renewcommand{\sectionmark}[1]{\markleft{\thesection.\ \ #1}}
	\ofoot{\pagemark}
	\cfoot{}

% clereref must be loaded after anything that changes the referencing system
\usepackage[capitalise, noabbrev]{cleveref} % these options keep the word "Figure"
\creflabelformat{equation}{#2#1#3}

% make the code float work with cleverref
\crefname{codecell}{code}{codes}
\Crefname{codecell}{code}{codes}
% make the text float work with cleverref
\crefname{textcell}{text}{texts}
\Crefname{textcell}{text}{texts}
% make the text float work with cleverref
\crefname{errorcell}{error}{errors}
\Crefname{errorcell}{error}{errors}

	\begin{document}

		\begin{titlepage}
	\begin{center}
		\includegraphics[width=0.3\textwidth]{template_files/princeton_logo.pdf}
	\end{center}

	\begin{center}

	\vspace*{1cm}

	\Huge\textbf{PsiPyPublish: An IPyPublish Template for Psychological Research}

	\vspace{0.8cm}

	\begin{center}\LARGE{Stefan Uddenberg}
	\end{center}
		  \LARGE{Princeton University}\\

	\begin{flushleft} \large\vspace{1.2cm}
		\emph{Running head:} PsiPyPublish
	\end{flushleft}

	\begin{flushleft} \large\emph{Address for correspondence:} \\
			Stefan Uddenberg\\
			Peretsman Scully Hall 322\\
			Princeton University\\
			Princeton, NJ 08540\\
			stefanu@princeton.edu\\
	\end{flushleft}

	\begin{flushleft} \large\emph{Word count:} XXX (Main text + abstract)
	\end{flushleft}

	\begin{flushleft} \large\emph{Version:} Template file -- not for submission
	\end{flushleft}

	\vfill

	\begin{minipage}{0.8\textwidth}
	\begin{center}
	\end{center}
	\end{minipage}

	\today

	\end{center}
	\end{titlepage}

		\begingroup
    \let\cleardoublepage\relax
    \let\clearpage\relax% remove header and footer from Table of Contents
	\thispagestyle{empty}
    \tableofcontents
    \endgroup

\section{Introduction}\label{introduction}

This is a template for an APA-style
\href{https://github.com/chrisjsewell/ipypublish}{iPyPublish}
manuscript. Feel free to check out the documentation and examples at
that link; it's all very good. There you can find information on how to
embed figures, code, tables, and more. References are managed using
\href{https://www.zotero.org/}{Zotero} in concert with
\href{https://github.com/retorquere/zotero-better-bibtex/}{Better
BibTex}. For now, you're going to want to edit the notebook's metadata
in order to change what appears on the title page. In addition, the
metadata includes \texttt{jupytext} configuration, so that you can
automatically generate markdown and py:percent versions of this notebook
automatically on saving -\/- assuming you have \texttt{jupytext}
installed and correctly configured, that is!

\subsection{Configuration}\label{configuration}

My working configuration files for Jupyter (with Jupytext) and
iPyPublish can be found in this repository. Naturally, you will need to
replace your computer's original versions of these files with the new
ones included here. For example, if using Anaconda, your iPyPublish
installation can be found at
\texttt{your\_environment\_name/Lib/site-packages/ipypublish}.

\begin{itemize}
\tightlist
\item
  \texttt{biblio\_natbib.py}, \texttt{doc\_article.py}, and
  \texttt{front\_pages.py} all live in
  \texttt{your\_environment\_name\textbackslash{}Lib\textbackslash{}site-packages\textbackslash{}ipypublish\textbackslash{}latex\textbackslash{}ipypublish}
\item
  \texttt{ipynb\_latex\_setup.py} lives in
  \texttt{your\_environment\_name\textbackslash{}Lib\textbackslash{}site-packages\textbackslash{}ipypublish\textbackslash{}scripts}
\item
  \texttt{latex\_ipypublish\_main.py} lives in
  \texttt{your\_environment\_name\textbackslash{}Lib\textbackslash{}site-packages\textbackslash{}ipypublish\textbackslash{}export\_plugins}
\end{itemize}

\subsection{Caveats}\label{caveats}

Since the creation of this template, \texttt{ipypublish} has been
upgraded to version 0.9.0. This template was designed to work with
version 0.6.7, and suits my needs; as such, it may take some time before
I update this guide to deal with the latest version --- if any changes
are even needed at all, since I haven't had a chance to try out the
latest edition.

\section{Notes}\label{notes}

\subsection{Production}\label{production}

Produce a notebook in the terminal with the command
\texttt{nbpublish\ -pdf\ -pbug\ file\_name.ipynb} \footnote{\hphantom{}
  Technically \texttt{-pbug} is optional so you can see verbose output,
  but nbpublish seems to work more reliably with this option enabled.}.
Outputs to \texttt{converted} folder at the \texttt{.ipynb} file's
location.\footnote{\hphantom{} \texttt{nbpublish} requires a lot of
  different technologies to work together. As such, if a build fails,
  simply try running the same command once more to see if that fixes the
  issue before moving on to more intense debugging.}

\subsection{Markdown}\label{markdown}

\begin{itemize}
\item
  Headings and sub-headings (and so on) are made by prefacing text with
  \texttt{\#}. The more \texttt{\#}s, the greater the level of heading.
\item
  Unordered lists are made by prefacing text with a "-".

  \begin{enumerate}
  \def\labelenumi{\arabic{enumi}.}
  \tightlist
  \item
    Numbered lists start with a number and dot.
  \item
    Create sublists via tabbed indentation.
  \end{enumerate}
\item
  Footnote links are made with \texttt{{[}\^{}X{]}} (where \texttt{X} is
  some number). Footnote content is placed below with
  \texttt{{[}\^{}X{]}:\ Content\ goes\ here}. Here's an
  example.\footnote{\hphantom{} Footnote content goes here!}

  \begin{itemize}
  \tightlist
  \item
    Correct formatting only appears after running \texttt{nbpublish}.
  \end{itemize}
\item
  \href{https://google.com}{Links} can be generated with the following
  syntax: \texttt{{[}link{]}(http://www.website.com)}
\item
  \texttt{Code} can be placed between backticks (the character to the
  left of the \texttt{1} key at the top of your keyboard).

  \begin{itemize}
  \tightlist
  \item
    Place it between 3 backticks (with an optional language name) and
    you get (syntax-highlighted) block code.\footnote{\hphantom{} Note,
      however, that one should not use this for displaying large chunks
      of code in an nbpublish PDF. Instead, see code cell
      \cref{code:publish} below for an example of how to place code in
      the PDF} \texttt{python\ print(foo)}
  \end{itemize}
\item
  \emph{Italic}, \textbf{bold}, and \textbf{\emph{bolded italic}} text
  can be created by sandwiching text between 1, 2, or 3 \texttt{*}s or
  \texttt{\_}s respectively.
\item
  \begin{quote}
  Blockquotes are made by prefacing text with \texttt{\textgreater{}} .
  \end{quote}
\end{itemize}

\todo[inline]{Get inline todos with Latex's "todo" command.}

\subsection{Citations and References}\label{citations-and-references}

\begin{itemize}
\tightlist
\item
  First, specify the \texttt{bibliography} entry in the notebook
  metadata to the correct bibliography file (Edit -\/-\textgreater{}
  Edit Notebook Metadata). \emph{Leave out the \texttt{.bib} extension
  from this file name!} It should look like \texttt{path/to/bibFileName}
  .

  \begin{itemize}
  \tightlist
  \item
    If nbpublish is having problems finding the \texttt{.bib} file, I
    have had success by placing a copy in the
    \texttt{converted/notebook\_name\_files/} directory, as well as
    placing the file in the same folder as the actual notebook. This
    makes set up for the notebook's bibliography metadata especially
    easy.
  \end{itemize}
\item
  Citations can be input with citation keys and standard latex commands
  (e.g., \texttt{\textbackslash{}cite\{citationKey\}}).
\item
  I've had success with citation keys generated via Zotero Better
  BibTex, like so \citep{uddenbergTelefaceSerialReproduction2018}. Note
  that you won't see the final formatted output until you run
  \texttt{nbpublish}.
\item
  See a \href{http://merkel.texture.rocks/Latex/natbib.php}{cheat sheet
  of valid cite commands here}.
\end{itemize}

\subsection{Figures}\label{figures}

\begin{itemize}
\item
  Figures can be displayed with commands like
  \texttt{display(SVG("filename.svg"))} or
  \texttt{Image(\textquotesingle{}filename.jpg\textquotesingle{},\ height=400)}.
\item
  Edit the cell's metadata to change the figure caption, placement,
  size, et al. (View -\/-\textgreater{} Cell Toolbar -\/-\textgreater{}
  Edit Metadata -\/-\textgreater{} Click on "Edit Metadata" above cell.)
\item
  Figures can be referenced via
  \texttt{\textbackslash{}cref\{fig:figNameFromMetadata\}}. For example:
  \cref{fig:example}
\item
  Here's the metadata for the figure below (so you don't have to inspect
  it yourself):

\begin{Shaded}
\begin{Highlighting}[]
\FunctionTok{\{}
  \DataTypeTok{"ipub"}\FunctionTok{:} \FunctionTok{\{}
\DataTypeTok{"figure"}\FunctionTok{:} \FunctionTok{\{}
  \DataTypeTok{"caption"}\FunctionTok{:} \StringTok{"An example beeswarm plot of the built-in `tips` dataset found in seaborn."}\FunctionTok{,}
  \DataTypeTok{"height"}\FunctionTok{:} \FloatTok{0.4}\FunctionTok{,}
  \DataTypeTok{"label"}\FunctionTok{:} \StringTok{"fig:example"}\FunctionTok{,}
  \DataTypeTok{"placement"}\FunctionTok{:} \StringTok{"H"}\FunctionTok{,}
  \DataTypeTok{"widefigure"}\FunctionTok{:} \KeywordTok{false}
\FunctionTok{\}}
  \FunctionTok{\}}
\FunctionTok{\}}
\end{Highlighting}
\end{Shaded}
\end{itemize}

\subsection{Displaying Code}\label{displaying-code}

Displaying the code in a code block, as in \cref{code:example_list_comp}
can be accomplished by editing the metadata:

\begin{Shaded}
\begin{Highlighting}[]
\FunctionTok{\{}
\DataTypeTok{"ipub"}\FunctionTok{:} \FunctionTok{\{}
  \DataTypeTok{"code"}\FunctionTok{:} \FunctionTok{\{}
  \DataTypeTok{"format"} \FunctionTok{:} \FunctionTok{\{\},}
    \DataTypeTok{"asfloat"}\FunctionTok{:} \KeywordTok{true}\FunctionTok{,}
    \DataTypeTok{"caption"}\FunctionTok{:} \StringTok{"Example list comprehension."}\FunctionTok{,}
    \DataTypeTok{"label"}\FunctionTok{:} \StringTok{"code:example_list_comp"}\FunctionTok{,}
    \DataTypeTok{"widefigure"}\FunctionTok{:} \KeywordTok{false}\FunctionTok{,}
    \DataTypeTok{"placement"}\FunctionTok{:} \StringTok{"H"}
    \FunctionTok{\}}
  \FunctionTok{\}}
\FunctionTok{\}}
\end{Highlighting}
\end{Shaded}

\begin{codecell}[H]

    \caption{Example list comprehension}\label{code:example_list_comp}\begin{lstlisting}[language=Python,numbers=left,xleftmargin=20pt,xrightmargin=5pt,belowskip=5pt,aboveskip=5pt]
colors = ["red", "green", "blue"]
for color in colors:
    print(color)
\end{lstlisting}\end{codecell}

\subsection{Templating --- Pass Variables into
Markdown}\label{templating-pass-variables-into-markdown}

\begin{itemize}
\tightlist
\item
  Using the
  \href{https://jupyter-contrib-nbextensions.readthedocs.io/en/latest/nbextensions/python-markdown/readme.html}{Python
  Markdown Extension}, you can pipe valid Python code into markdown
  cells directly by sandwiching it between two curly braces: E.g., 2 + 2
  = 4
\item
  Note that the notebook needs to be \texttt{Trusted}; look to the top
  right to see if it is and simply click on \texttt{Not\ Trusted} to
  change that.
\end{itemize}

\subsection{Latex}\label{latex}

\begin{itemize}
\tightlist
\item
  Execute arbitrary latex by sandwiching it between dollar signs:
  \(a = b+c\)
\item
  Alternatively, use \texttt{Latex()} command from \texttt{ipypublish}
  within a code cell.
\item
  Latex's \texttt{hphantom} command is useful when you just want a
  little more horizontal space between items.
\end{itemize}

\subsection{Terminal commands}\label{terminal-commands}

\begin{itemize}
\tightlist
\item
  Execute terminal commands in Jupyter by prefacing code with \texttt{!}
  .
\item
  For example, you can export this notebook with the following code cell
  (uncommented, of course):
\end{itemize}

\begin{codecell}[H]

    \caption{How to publish this notebook, from within the notebook itself!}\label{code:publish}\begin{lstlisting}[language=Python,numbers=left,xleftmargin=20pt,xrightmargin=5pt,belowskip=5pt,aboveskip=5pt]
# !nbpublish -pdf -pbug template.ipynb
\end{lstlisting}\end{codecell}

\subsection{Troubleshooting}\label{troubleshooting}

\subsubsection{Jupytext}\label{jupytext}

\begin{itemize}
\tightlist
\item
  If saving to \texttt{Rmd} format, beware using single quotes within
  figure caption metadata, since R Markdown uses single quotes for
  metadata and not double quotes, which creates issues when you want to
  include apostrophes.
\item
  If you are encountering problems opening or even "trusting" a notebook
  that was previously working fine, simply delete the other non-ipynb
  representations of the notebook. I encounter this issue most often
  when synchronizing notebook files via Dropbox.
\end{itemize}

% sort citations by order of first appearance
\bibliographystyle{apacite} % formerly unsrtnat with natbib
\setlength\bibsep{\baselineskip} % increase space between refs
\bibliography{template_files/ExampleBib}

	\end{document}

